% =============================================================================
%  vsim_nomenclature_review.tex
%  Nomenclature Inconsistency Review
%  Bridge docs cross-checked against VSIM_REFERENCE.md (v5.0.2)
%  REV-2026-01
% =============================================================================
\documentclass[11pt,a4paper]{article}
\usepackage{vsim_common}
\usepackage{soul}       % for \st strikethrough

\hypersetup{pdftitle={Nomenclature Review REV-2026-01},pdfsubject={VSIM Internal}}

% Severity pills
\definecolor{critbg}{RGB}{253,232,232}
\definecolor{critfg}{RGB}{120,10,10}
\definecolor{sigbg}{RGB}{255,248,230}
\definecolor{sigfg}{RGB}{120,80,0}
\definecolor{minbg}{RGB}{244,244,244}
\definecolor{minfg}{RGB}{60,60,60}

\newcommand{\scrit}{\spill{critbg}{critfg}{CRITICAL}}
\newcommand{\ssig}{\spill{sigbg}{sigfg}{SIGNIFICANT}}
\newcommand{\sminor}{\spill{minbg}{minfg}{MINOR}}

% Finding box
\newcounter{finding}
\newcommand{\finding}[3]{%   #1=severity pill  #2=short title  #3=body
  \stepcounter{finding}%
  \vspace{8pt}%
  \begin{tcolorbox}[colback=white,colframe=black!20,arc=0pt,boxrule=0.5pt,
    boxsep=4pt,left=8pt,right=8pt,top=6pt,bottom=6pt,breakable]
    \noindent{\sffamily\bfseries\small F\thefinding\quad #2}\hfill #1\\[5pt]
    {\small #3}
  \end{tcolorbox}%
}

% Inline wrong/right
\newcommand{\wrong}[1]{\textcolor{failtext}{\ttfamily\small #1}}
\newcommand{\correct}[1]{\textcolor{passtext}{\ttfamily\small #1}}

\begin{document}
\thispagestyle{empty}

\begin{center}
  \eyebrow{VSIM Internal \textperiodcentered\ Nomenclature Review}\\[8pt]
  {\LARGE\sffamily\bfseries Nomenclature Inconsistency Review}\\[6pt]
  {\normalsize\sffamily
    Bridge Docs Cross-checked Against \texttt{VSIM\_REFERENCE.md} v5.0.2}
\end{center}

\vspace{10pt}
\noindent
\begin{tabular}{@{}L{2.8cm} L{9cm}@{}}
  \code{Doc ID}     & REV-2026-01 \\
  \code{Scope}      & \texttt{vsim\_bridge\_doc.tex}, \texttt{vsim\_code\_request.tex},
                      \texttt{vsim\_bridge\_audit.tex} \\
  \code{Reference}  & \texttt{VSIM\_REFERENCE.md} (v5.0.2) \\
  \code{Date}       & 2026-06-01 \\
  \code{Status}     & Open --- findings require author action \\
\end{tabular}

\vspace{10pt}
\begin{infobox}
  \small\sffamily
  \textbf{Severity guide:}\quad
  \scrit\ = conflicts directly with reference; breaks existing parser or scripts\\[3pt]
  \ssig\ = struct/TOML mismatch or naming drift; will cause silent bugs\\[3pt]
  \sminor\ = style inconsistency or ambiguity; no runtime impact
\end{infobox}

\noindent\rule{\linewidth}{2pt}\vspace{4pt}
\tableofcontents
\newpage

% =============================================================================
\section{Critical Findings}
% =============================================================================

\finding{\scrit}{%
  \texttt{[observe].track} --- wrong key name
}{%
  \textbf{Location:} \texttt{vsim\_bridge\_doc.tex} \S\,9 demo script (line 380);
  \texttt{vsim\_code\_request.tex} CR-08 task description.\\[4pt]
  \textbf{Used:} \wrong{track = ["contact", "residence", ...]}\\
  \textbf{Reference says:} \correct{metrics = ["energy\_map", "rdf", ...]}\\[4pt]
  \texttt{VSIM\_REFERENCE.md} \S\,\texttt{[observe]} is unambiguous ---
  the key is \texttt{metrics}, not \texttt{track}.
  Every existing suite script (\texttt{suite\_gas\_plasma}, \texttt{suite\_metal\_steam})
  uses \texttt{metrics}. A parser will either reject \texttt{track} or silently ignore it,
  meaning the \texttt{identity\_entropy} observable would never activate.
  CR-08 propagates the same wrong key in its task description.\\[4pt]
  \textbf{Fix:} Replace all occurrences of \wrong{track} with \correct{metrics}
  in the demo script and in CR-08.
}

\finding{\scrit}{%
  \texttt{[system]} block conflicts with chemistry legacy alias
}{%
  \textbf{Location:} \texttt{vsim\_bridge\_doc.tex} \S\,9 demo script (lines 357--359).\\[4pt]
  \textbf{Used:}
  \begin{databox}
  {\ttfamily\small [system]\quad mode = "surface\_interaction"\quad space = "3d"}
  \end{databox}
  \textbf{Reference says:} \texttt{[system]} is the \emph{legacy alias}
  for \texttt{[chemistry]} / \texttt{[chemistry].chemistry}. It is already a
  parsed section with completely different semantics (ambient reaction rules, heat gate,
  scoring thresholds). Using \texttt{[system]} with \texttt{mode} and \texttt{space} keys
  will be parsed as chemistry configuration and produce undefined behaviour.\\[4pt]
  \textbf{Fix:} Rename to a non-conflicting block name. Candidates:
  \correct{[run.surface]}, \correct{[surface\_mode]}, or promote these fields into
  \texttt{[run]} (which already has \texttt{mode = "md"}).
  The \texttt{[[surface]]} array-of-tables block below it should remain.
}

\finding{\scrit}{%
  Non-standard metric names in \texttt{[observe].metrics}
}{%
  \textbf{Location:} \texttt{vsim\_bridge\_doc.tex} \S\,9 demo script.\\[4pt]
  \textbf{Used:} \wrong{"contact"}, \wrong{"residence"}, \wrong{"embedding"},
  \wrong{"diffusion\_axis"}\\[4pt]
  None of these appear in the \texttt{VSIM\_REFERENCE.md} observe metrics table.
  The closest match is \texttt{"diffusion\_proxy"} (reference) vs
  \wrong{"diffusion\_axis"} (bridge doc) --- different names for what may be the same concept.
  \wrong{"contact"}, \wrong{"residence"}, and \wrong{"embedding"} appear to be
  surface-interaction-specific observables not yet registered in the reference.\\[4pt]
  \textbf{Fix:} Either register these names in the observe metrics table in
  \texttt{VSIM\_REFERENCE.md} (if they are real planned metrics), or remove them
  from the demo script. Replace \wrong{"diffusion\_axis"} with
  \correct{"diffusion\_proxy"} if the intent is the existing metric.
}

\finding{\scrit}{%
  \texttt{[particles.projectile]} is undocumented; conflicts with \texttt{[[simulation.molecule]]}
}{%
  \textbf{Location:} \texttt{vsim\_bridge\_doc.tex} \S\,9 demo script (lines 366--370).\\[4pt]
  \textbf{Used:}
  \begin{databox}
  {\ttfamily\small
  [particles.projectile]\\
  count = 128\quad type = "generic\_bead"\quad velocity = [0.0, 0.0, -1.5]\quad height = 15.0}
  \end{databox}
  The canonical way to declare particle species in VSIM is
  \correct{[[simulation.molecule]]} (with \texttt{formula}, \texttt{count},
  \texttt{temperature}, \texttt{lattice}). The bridge doc introduces
  \texttt{[particles.projectile]} with entirely different field names (\texttt{type}
  instead of \texttt{formula}; \texttt{height} and \texttt{velocity} as top-level
  keys instead of via \texttt{[[override.particle]]}). The parser will not recognise
  this block.\\[4pt]
  \textbf{Fix:} Replace with \correct{[[simulation.molecule]]} + initial velocity
  via \correct{[[override.particle]]} (which IS in the reference), or register
  \texttt{[particles.*]} as a new section type in the reference and parser.
}

% =============================================================================
\section{Significant Findings}
% =============================================================================

\finding{\ssig}{%
  \texttt{surface\_state} field in TOML missing from \texttt{FormationContext} C++ struct
}{%
  \textbf{Location:} \S\,9 demo script vs \S\,10 C++ header.\\[4pt]
  Demo script \texttt{[formation.initial\_state]} contains:
  \wrong{surface\_state = "roughened"}\\
  The \texttt{FormationContext} struct in \S\,10 has fields:
  \code{method}, \code{history}, \code{formation\_time}, \code{temperature\_K},
  \code{pressure\_Pa}, \code{phase}, \code{defect\_state}, \code{grain\_state},
  \code{density} --- \textbf{no} \texttt{surface\_state}.\\[4pt]
  CR-01 in the code request also lists the same nine fields with no \texttt{surface\_state}.
  The field would be silently dropped by the parser. Either add it to the struct and
  the CR-01 field list, or remove it from the TOML example.
}

\finding{\ssig}{%
  TOML field name \texttt{temperature} vs C++ struct field \texttt{temperature\_K}
}{%
  \textbf{Location:} \S\,4 minimum formation block vs \S\,10 C++ struct.\\[4pt]
  \S\,4 TOML block: \wrong{temperature = 300.0} (no unit suffix)\\
  \S\,10 C++ struct: \correct{double temperature\_K = 300.0;} (with \texttt{\_K} suffix)\\[4pt]
  The \S\,9 demo script correctly uses \texttt{temperature = 300.0} in
  \texttt{[formation.initial\_state]}, which maps to the C++ field \texttt{temperature\_K}.
  This TOML-to-struct name mismatch must be documented in the parser (similar to how
  \texttt{[run]} accepts both \texttt{temperature} and \texttt{temperature\_K}).
  The \S\,4 minimum block should note the unit convention or use the suffixed name
  for clarity.
}

\finding{\ssig}{%
  TOML field \texttt{pressure} vs C++ struct field \texttt{pressure\_Pa}
}{%
  \textbf{Location:} \S\,4 minimum formation block vs \S\,10 C++ struct.\\[4pt]
  Same issue as F6: the TOML uses \wrong{pressure = 101325.0} while the struct
  declares \correct{double pressure\_Pa = 101325.0;}. Note also that \texttt{[run]}
  uses \texttt{pressure\_GPa} (different unit), creating three pressure representations
  across the system. The formation block implicitly uses Pascals but the field name
  does not say so. The unit should be either embedded in the TOML key
  (\correct{pressure\_Pa}) or documented in a comment.
}

\finding{\ssig}{%
  Column header abbreviations in \texttt{.ie} data row do not match field names
}{%
  \textbf{Location:} \texttt{vsim\_bridge\_doc.tex} \S\,6 data row example.\\[4pt]
  The example header uses shortened names that do not match the formal field names
  written by \texttt{IEWriter::write\_header()} in \S\,10:

  \vspace{4pt}
  \begin{center}\small
  \begin{tabular}{@{}ll@{}}
    \rowcolor{skybg}\textbf{Data row (abbreviated)} & \textbf{Formal field name} \\
    \toprule
    \wrong{hidden\_act}  & \correct{hidden\_activity} \\
    \wrong{proj\_loss}   & \correct{projection\_loss} \\
    \wrong{recov}        & \correct{recoverability} \\
    \wrong{coal\_loss}   & \correct{coalescence\_loss} \\
    \wrong{lin}          & \correct{lineage\_depth} \\
    \wrong{state\_q}     & \correct{state\_quota} \\
    \wrong{entr\_flux}   & \correct{entropy\_flux} \\
    \wrong{form\_mem}    & \correct{formation\_memory} \\
    \bottomrule
  \end{tabular}
  \end{center}
  \vspace{4pt}
  The actual \texttt{.ie} file written by \texttt{IEWriter} uses full names
  (confirmed in \S\,10 \texttt{write\_header()}). The example would mislead anyone
  using it to write a parser for the \texttt{.ie} format.\\[4pt]
  \textbf{Fix:} Use full field names in the data row example, or note explicitly
  that the header row is abbreviated for display purposes only.
}

\finding{\ssig}{%
  \texttt{[seed].formation} vs new \texttt{[formation]} block --- namespace collision risk
}{%
  \textbf{Location:} \texttt{VSIM\_REFERENCE.md} \S\,\texttt{[seed]} vs
  bridge doc \S\,4.\\[4pt]
  The existing \texttt{[seed]} section has a field \texttt{formation = 7306}
  (a 64-bit RNG seed component). The bridge doc introduces a new top-level block
  \texttt{[formation]} with entirely different semantics (material history context).
  The name proximity is a parser confusion risk: a naive reader of
  \code{seed.formation = 7306} may conflate this with the new \texttt{[formation]}
  block. The suite scripts already use both: \code{[seed] formation = 8106} alongside
  the intended new \code{[formation]} block. The reference should add a disambiguation
  note, and the bridge doc should acknowledge this naming overlap.
}

\finding{\ssig}{%
  \texttt{[[surface]]} block is undocumented in \texttt{VSIM\_REFERENCE.md}
}{%
  \textbf{Location:} \texttt{vsim\_bridge\_doc.tex} \S\,9 demo script.\\[4pt]
  The demo uses a \texttt{[[surface]]} block (array-of-tables) with fields
  \code{type = "graphite\_layer"}, \code{roughness = 0.05},
  \code{warm\_minimum\_beads = 100}.\\[4pt]
  \texttt{[[surface]]} (array-of-tables), its fields \texttt{roughness},
  \texttt{warm\_minimum\_beads}, and the type value \texttt{"graphite\_layer"}
  do not appear anywhere in \texttt{VSIM\_REFERENCE.md}. Either these are planned
  future additions that need registering, or the block name conflicts with something
  not yet surfaced in the reference.
}

% =============================================================================
\section{Minor Findings}
% =============================================================================

\finding{\sminor}{%
  \texttt{defect\_state} (field) vs ``defect seed'' (prose) --- semantic ambiguity
}{%
  \textbf{Location:} \S\,8 formation descriptors list vs \S\,10 C++ struct.\\[4pt]
  The formation descriptors table in \S\,8 lists ``defect seed'' as a descriptor type.
  The \texttt{FormationContext} struct field is \texttt{defect\_state}. These are
  different concepts: a \emph{seed} is a numerical RNG input for defect placement;
  a \emph{state} is the current defect configuration. The VSIM\_REFERENCE
  \texttt{[seed]} section already has a \texttt{defect} seed field (\texttt{defect = 7303}).
  Recommend changing ``defect seed'' in the \S\,8 prose to ``defect state'' to match
  the struct, or renaming the struct field to \texttt{defect\_seed} if the intent is
  to store a seed value.
}

\finding{\sminor}{%
  \texttt{formation\_time} units unspecified
}{%
  \textbf{Location:} \S\,4 minimum formation block.\\[4pt]
  \code{formation\_time = 0.0} has no documented unit. All time-related fields in
  the reference use explicit unit suffixes (e.g., \texttt{dt\_fs} in femtoseconds,
  \texttt{step\_delay\_ms} in milliseconds). Recommend either renaming to
  \correct{formation\_time\_fs} or adding a comment in the TOML schema.
}

\finding{\sminor}{%
  \texttt{[outputs.identity\_entropy]} not following \texttt{[export.*]} convention
}{%
  \textbf{Location:} All bridge docs.\\[4pt]
  Already flagged as DEC-01 in the code request, but worth stating explicitly
  as a nomenclature finding. All existing output configuration uses the
  \texttt{[export]} and \texttt{[export.visual]} namespace. The bridge doc
  introduces \texttt{[outputs.*]}, a new top-level namespace that does not exist
  elsewhere in the reference. This creates two parallel output configuration
  systems with no clear rule for which one owns what. DEC-01 must be resolved
  and the decision recorded in the reference.
}

\finding{\sminor}{%
  \texttt{vsim\_bridge\_audit.tex} audit tables use ``Wired'' where reference uses \checkmark
}{%
  \textbf{Location:} \texttt{vsim\_bridge\_audit.tex} \S\,3 Runtime Wire Audit.\\[4pt]
  The audit doc uses the pill \swired\ for existing functions. The
  \texttt{VSIM\_REFERENCE.md} uses the symbol legend: \checkmark\ (parsed and wired),
  \texttt{*} (parsed, not wired), \texttt{x} (not parsed). The audit doc should note
  that its status pills do not correspond to the reference symbols, or adopt the same
  convention to reduce cross-document confusion.
}

\finding{\sminor}{%
  \texttt{[run].temperature\_K} vs \texttt{[formation.initial\_state].temperature}
}{%
  \textbf{Location:} \S\,9 demo script.\\[4pt]
  The \S\,9 demo uses \wrong{temperature\_K = 300.0} in \texttt{[run]} and
  \wrong{temperature = 300.0} in \texttt{[formation.initial\_state]}.
  These are different field names for the same physical quantity (temperature in Kelvin).
  The reference notes that \texttt{[run]} accepts both \texttt{temperature} and
  \texttt{temperature\_K} as aliases, so both are technically valid, but mixing
  them within a single script increases read-time confusion. Recommend using the
  suffixed form (\correct{temperature\_K}) consistently across all blocks in new scripts.
}

\finding{\sminor}{%
  \texttt{C++ IERecord.state\_quota} is \texttt{uint32\_t} but \texttt{lineage\_depth} range may exceed 32-bit on deep coalescence chains
}{%
  \textbf{Location:} \S\,10 C++ header.\\[4pt]
  \code{uint32\_t lineage\_depth = 0;} and \code{uint32\_t state\_quota = 1;}
  are both 32-bit unsigned. For \texttt{state\_quota} ($Q_n$, particle count),
  \texttt{uint32\_t} caps at $\approx 4\times10^9$, which is sufficient. For
  \texttt{lineage\_depth}, the cap is also fine in practice. However, the frame
  and \texttt{identity\_id} fields are \texttt{uint64\_t}. The mixed 32/64-bit
  usage should be documented as a deliberate choice in the header, or both counters
  widened to \texttt{uint64\_t} for consistency.
}

% =============================================================================
\section{Summary Table}
% =============================================================================

\vspace{6pt}
\small
\begin{longtable}{@{} C{0.8cm} L{1.5cm} L{7.5cm} L{3.2cm} @{}}
  \rowcolor{skybg}
  \textbf{F\#} & \textbf{Severity} & \textbf{Title} & \textbf{Location} \\
  \toprule
  \endhead\bottomrule\endfoot

  F1  & \scrit   & \code{[observe].track} wrong key name             & bridge \S\,9, CR-08 \\
  F2  & \scrit   & \code{[system]} conflicts with chemistry alias     & bridge \S\,9 \\
  F3  & \scrit   & Non-standard observe metric names                  & bridge \S\,9 \\
  F4  & \scrit   & \code{[particles.projectile]} undocumented         & bridge \S\,9 \\
  F5  & \ssig    & \code{surface\_state} in TOML, absent from struct  & bridge \S\,9, \S\,10 \\
  F6  & \ssig    & TOML \code{temperature} vs struct \code{temperature\_K} & bridge \S\,4, \S\,10 \\
  F7  & \ssig    & TOML \code{pressure} vs struct \code{pressure\_Pa} & bridge \S\,4, \S\,10 \\
  F8  & \ssig    & Column header abbreviations don't match field names & bridge \S\,6 \\
  F9  & \ssig    & \code{[seed].formation} vs \code{[formation]} collision risk & ref + bridge \S\,4 \\
  F10 & \ssig    & \code{[[surface]]} block undocumented              & bridge \S\,9 \\
  F11 & \sminor  & \code{defect\_state} vs ``defect seed'' prose      & bridge \S\,8, \S\,10 \\
  F12 & \sminor  & \code{formation\_time} units unspecified           & bridge \S\,4 \\
  F13 & \sminor  & \code{[outputs.*]} vs \code{[export.*]} namespace  & all bridge docs \\
  F14 & \sminor  & Audit status pills vs reference symbol legend      & audit \S\,3 \\
  F15 & \sminor  & Mixed \code{temperature} / \code{temperature\_K} in one script & bridge \S\,9 \\
  F16 & \sminor  & Mixed \code{uint32\_t} / \code{uint64\_t} widths in IERecord  & bridge \S\,10 \\

\end{longtable}
\normalsize

\vspace{8pt}
\begin{warnbox}
  \textbf{Recommended action order:} Fix F1 (wrong observe key) and F2 (system block conflict)
  immediately --- both would cause a real script to silently misbehave or error at parse time.
  F3 and F4 (undocumented blocks/metrics) should be resolved before the bridge doc \S\,9
  script is used as a canonical example. F5--F10 should be resolved before CR-01 is closed.
\end{warnbox}

% =============================================================================
\vspace{12pt}
\noindent\rule{\linewidth}{0.5pt}
{\small\sffamily\textcolor{gray}{REV-2026-01 \textperiodcentered\
  VSIM Internal \textperiodcentered\
  Ref: VSIM\_REFERENCE.md v5.0.2 \textperiodcentered\
  vsim\_bridge\_doc.tex \textperiodcentered\ vsim\_code\_request.tex}}
\end{document}
